\chapter{Verifizierung}\label{chap:Verifizierung}
Die Einhaltung der in Polarion\footnote{\url{https://plm.sw.siemens.com/de-DE/polarion/} --\,zuletzt am 21.~Februar~2025 erfolgreich aufgerufen.} dokumentierten Spezifikationen ist ein zentraler Bestandteil der Entwicklung der \gls{por}-Schaltung. Diese Spezifikationen definieren die funktionskritischen Parameter der Schaltung, wie beispielsweise die Einschaltverzögerung, die Reset-Aktivierungsschwelle und die Reaktionszeit bei Spannungsabfall. Da der \gls{por} eine essenzielle Rolle für das fehlerfreie Hochfahren und Zurücksetzen des gesamten Systems spielt, müssen diese Parameter innerhalb der definierten Grenzen liegen, um eine zuverlässige Initialisierung des Chips sicherzustellen.

\section{Spezifikationen}
Nachdem die grundlegende Funktionsweise der einzelnen Funktionsblöcke erfolgreich simuliert und validiert wurde, erfolgte eine gezielte Anpassung der Bauelemente, um die spezifizierten Anforderungen zu erfüllen. Die Schaltung wurde mit der Cadence\footnote{\url{https://www.cadence.com/en_US/home.html}  --\,zuletzt am 21.~Februar~2025 erfolgreich aufgerufen.}-Virtuoso-Design-Suite entworfen und getestet. Der Fokus lag dabei auf der korrekten Implementierung der verzögerten Aktivierung sowie der schnellen Reaktion des Reset-Signals beim Abschalten der Versorgungsspannung und dem Einstellen der Schwell- und Reset-Spannung des Komparators. 

Zunächst erfolgte die Simulation unter nominalen Bedingungen, um sicherzustellen, dass die Funktionsblöcke sich in Kombination wie zu erwarten verhalten. Zu Beginn wurden die Schaltschwellen des Komparators mit Hilfe der Widerstände eingestellt, um sichergehen zu können, dass sich die Reset-Spannung ($V_L$ siehe \autoref{fig:Reset-Spannung}) unter Schwellspannung ($V_H$ siehe \autoref{fig:Schwellspannung})befindet. Anschließend wurde besonders die Lade- und Entladezeit des Kondensators für die Reset-Freigabe analysiert, damit sich im System keine undefinierten Zustände einstellen, wenn es zu Spannungsschwankungen kommt. Durch Anpassungen und eine erneute Simulation konnte die Grundfunktionalität der \gls{por}-Schaltung verifiziert werden. Um die vorgegeben Spezifikationen bestmöglich zu erreichen, wurden gezielte Anpassungen an der Dimensionierung der Transistoren vorgenommen. Nachdem sich die Werte der Schaltung in dem Bereich der Spezifikationen befanden, wurde die nominale Corner-Simulation bei 27$^{\circ}$C, um verschiedene Prozesstypen sowie die maximalen Temperaturen erweitert. 

\section{Corner-Simulation}
Die Entwicklung einer robusten \gls{por}-Schaltung erfordert eine detaillierte Analyse der Schaltungsperformance unter verschiedenen Betriebsbedingungen. In modernen Halbleiterprozessen unterliegen Transistoren und passive Bauelemente\footnote{Dazu zählen u.a. Widerstände und Kondensatoren} signifikanten Variationen, die durch Unterschiede in der Produktion, Umgebungstemperatur und Versorgungsspannung verursacht werden. Um die korrekte Funktion der Schaltung sicherzustellen, wurde zuerst eine Corner-Simulation durchgeführt. Diese Simulation ermöglicht es, die Auswirkungen von \gls{pvt} systematisch zu untersuchen und gegebenenfalls gezielte Entwurfsanpassungen vorzunehmen. Werden Prozess- und Temperaturvariationen unzureichend berücksichtigt, kann dies dazu führen, dass die \gls{por}-Schaltung unter bestimmten Bedingungen fehlerhaft arbeitet. Beispielsweise könnte die Schaltung in einem \enquote{schnellen} Prozess-Corner~(höhere Elektronenmobilität, geringer Widerstände) zu früh oder in einem \enquote{langsamen} Prozess-Corner~(geringere Mobilität, höhere Widerstände) zu spät auslösen. Beides könnte schwerwiegende Konsequenzen für das Gesamtsystem haben, da Teile der restlichen Schaltung möglicherweise in einem undefinierten Zustand starten. In der Cadence-Virtuoso-ADE-Explorer-Umgebung werden daher verschiede Prozessvariationen und Temperaturbedingungen systematisch getestet. Die Standard-Prozesscorners in modernen CMOS-Technologien lassen sich wie folgt klassifizieren: 
\begin{table}[h!]
    \centering
    \footnotesize
    \begin{tabular}{c|c|c|c|c}
        \textbf{\makecell{Abkürzung}} &
        \textbf{\makecell{Bedeutung}} &
        \textbf{\makecell{Verhalten\\ NMOS}} &
        \textbf{\makecell{Verhalten\\ PMOS}} &
        \textbf{\makecell{Verhalten passive\\ Bauelemente}} \\
        \midrule
        ttt & typical-typical-typical & Nominal & Nominal & Nominal \\
        sss & slow-slow-slow & Langsam & Langsam & Langsam \\
        fff & fast-fast-fast & Schnell & Schnell & Schnell \\
        sft & slow-fast-typical & Langsam & Schnell & Nominal \\
        fst & fast-slow-typical & Schnell & Langsam & Nominal \\
        sfs & slow-fast-slow & Langsam & Schnell & Langsam \\
        sff & slow-fast-fast & Langsam & Schnell & Schnell \\
        fss & fast-slow-slow & Schnell & Langsam & Langsam \\
        fsf & fast-slow-fast & Schnell & Langsam & Schnell \\
    \end{tabular}
    \caption{Transistor- und Bauelement-Konfigurationen der Corner-Simulation}
    \label{tab:transistor_configurations}
\end{table}

Bei dieser Simulation können verschiedene Probleme auftreten, die eine Anpassung des Entwurfs erforderlich machen. Die Schwellspannung von \gls{nmos}- und \gls{pmos}-Transistoren kann sich je nach Prozess-Corner erheblich unterscheiden. In langsamen Corners (sss) ist die Schwellspannung typischerweise höher, wodurch sich die Transistoren langsamer schalten. In schnellen Corners (fff) wird diese reduziert, sodass die Transistoren früher schalten. Dadurch verschiebt sich die Reset-Schwelle der \gls{por}-Schaltung, was eine fehlerhafte oder verzögerte Initialisierung des Systems zur Folge haben kann. Die Kapazität des Kondensators~$C$ kann aufgrund von Fertigungstoleranzen schwanken. In einem sss-Prozess kann die Kapazität größer sein als erwartet, was die Ladezeit verlängert. In einem fff-Prozess ist die Kapazität möglicherweise geringer, sodass die Reset-Verzögerung verkürzt wird. 

Ohmsche Widerstände und Kapazitäten sind stark temperaturabhängig. Bei niedrigen Temperaturen von -40$^{\circ}C$ kann der Widerstand ansteigen, wodurch eine Verzögerung in der Schaltung entstehen kann. Hohe Temperaturen während des Betriebs können zu erhöhten Leckströmen und einer ungewollten Entladung des Kondensators führen, wodurch das Reset zu früh freigegeben wird. Während Prozessen wie sft (slow-fast-typical) oder fst (fast-slow-typical) arbeiten \gls{nmos}- und \gls{pmos}-Transistoren mit unterschiedlicher Geschwindigkeit. Dies kann zu asymmetrischem Verhalten im Schaltungsentwurf führen. Bei \gls{por}-Schaltungen muss das Reset-Signal sauber und zuverlässig detektiert werden, weshalb auch diese Simulation kritisch betrachtet werden muss.

Die Simulationsergebnisse (siehe \autoref{tab:corner_old_summary}) zeigen, dass die \gls{por}-Schaltung nur für wenige Corner-Fälle den vorgegebenen Spezifikationen entspricht, was auch durch Anpassungen der Kanalbreite und Länge der einzelnen Transistoren nicht stabilisiert werden kann. Grund dafür ist, dass Ohmsche Widerstände in integrierten Schaltungen signifikanten Fertigungsstreuungen aufgrund von Schwankungen in der Dicke und Dotierung des Halbleitermaterials unterliegen. Aber vor allem der Widerstandswert ändert sich mit der Temperatur, da der \gls{tcr} von Wiederständen relativ hoch ist \cite{ADofIntegratedCircuits}. Dadurch sind die Schwankungen der Referenzspannung (siehe \autoref{fig:Vref_fluctuation_old}) mit $\approx$ 320mV zu groß, um die Schaltung über die Corner-Simulation den Spezifikationen anzupassen. Deswegen muss der Entwurf für die bestehende Generierung der Referenzspannung aus \autoref{fig:Vref_old} überarbeitet werden.    
{\begin{table}[H]
    \centering
    \small
    \begin{tabular}{c|l|c|c|c}
        {\rotatebox{90}{\parbox{1.8cm}{\textbf{Spannung}}}} &
        \textbf{\makecell{Werte}} &
        \textbf{\makecell{Spezifikation}} &
        \textbf{\makecell{Min. Wert}} &
        \textbf{\makecell{Max. Wert}} \\
        \midrule

        \multirow{5}{*}{\rotatebox{90}{5 V}} &
        Verzögerung steigende Flanke & 20 us $\leq$ x $\leq$ 130 us & 20.27 us & 117.9 us \\
        & Verzögerung fallende Flanke & $\leq$ 1 us & 37.99 ns & 123.8 ns \\
        & Schwellspannung steigend & 2.2 V $\leq$ x $\leq$ 2.6 V & \cellcolor{lightred}1.811 V & \cellcolor{lightred}2.946 V \\
        & Schwellspannung fallend & 2.1 V $\leq$ x $\leq$ 2.5 V & \cellcolor{lightred}1.743 V & \cellcolor{lightred}2.849 V \\
        & Hysterese & 50 mV $\leq$ x $\leq$ 120 mV & 67.85 mV & 98.69 mV \\
        \midrule
    \end{tabular}
    \caption{Ergebnisse Corner-Simulation Widerstands-basierte PoR-Schaltung}
    \label{tab:corner_old_summary}
\end{table}

\subsection{Anpassung Referenzspannung}
Im Gegensatz dazu bietet die auf Transistoren basierende Referenzspannung (siehe \autoref{fig:Vref&Vres}) eine bessere Kompensation gegenüber Prozess- und Temperaturabweichungen. MOS-Tran\-sistoren in speziellen Konfigurationen, beispielsweise in Spannungsteiler- oder Stromspiegelstrukturen, können die Referenzspannung deutlich stabilisieren, da sie durch einen gezielten Schaltungsentwurf eine geringere Sensitivität gegenüber Fertigungsstreuungen aufweisen \cite{DesignOfAnalogCMOSIntegratedCircuits}.
In dem neuen Entwurf wird das Gate des nativen NMOS-Transistors $N_1$ durch die Spannung $V_{gate}$ gesteuert. $V_{gate}$ wiederum wird durch den Spannungsteiler aus $R_0$ und $N_0$ generiert. $N_1$ und $N_2$ bilden eine spannungsgesteuerte Referenzquelle, die durch die Spannung von $V_{gate}$ auch über Prozess- und Temperaturschwankungen konstanter ist. Damit der Komparator eine stabile Referenzspannung für die Spannungsversorgung hat, wurde der Entwurf für die Generierung der Referenzspannung (siehe \autoref{fig:Vref&Vres}) angepasst.  
%\vspace{1.5\baselineskip}
\begin{figure}[H]
    \begin{center}
        \begin{circuitikz}[european]
            \ctikzset{tripoles/mos style/arrows}[line width=thick, european]
            % Widerstand
            \draw (-4,4) coordinate(resistor-begin) to[resistor, l={$R_0$}] (-4,1) coordinate(resistor-end);
                %(-5.5,3) -- (-4.5,3) node[above, midway]{$Vdd$};
            
            % NMOS Transistor
            \draw (-4,-0.5) node[nmos, anchor=drain] (nmosL) {$N_0$};

            % NMOS Transistoren R 
            \draw (-1.5,3.25) node[nmosd, anchor=drain] (nmosL1) {$N_1$}
                %(nmosL1.gate) -- ++(-0.5,0) |- (resistor-end);
                (resistor-end) -- ++(1,0) |- (nmosL1.gate);
            \draw (-1.5,-1) node[nmos, anchor=drain] (nmosL2) {$N_2$}
                (nmosL1.source) -- (nmosL2.drain) coordinate [pos=0.5](nmosMP) % Verbindung & Mittelpunkt zwischen nmos
                (nmosMP) node[circle, fill, inner sep=1.5pt]{}
                (nmosL2.gate) -- ++(-0.25,0) |- (nmosMP)
                node[above, xshift=-20pt]{$V_{ref}$};
                
            
            % Verbindung von Drain des NMOS mit Widerstand
            \draw (resistor-end) -- (nmosL.drain);
            \draw (nmosL.gate) -- ++(-0.25,0) |- (resistor-end)
                node[circle, fill, inner sep=1.5pt]{}
                node[above, midway]{$V_{gate}$};
            %
            %*************************************************************************
            %
            % Widerstand R
            \draw (0,4.5) coordinate (resistor-begin1) to[resistor, l={$R_1$}] (0,2) coordinate (resistor-end1);
            \draw (0,2) coordinate (resistor-begin2) to[resistor, l={$R_2$}] (0,-1) coordinate (resistor-end2);
            \draw (0,-1) coordinate (resistor-begin3) to[resistor, l={$R_3$}] (0,-3.5) coordinate (resistor-end3);

            % NMOSR 
            \draw (3,2) node[nmos, rotate=270] (nmosR1) {}
                node at (nmosR1.center) [anchor=center, yshift=-10pt] {$N_3$}
                (nmosR1.source) -- (resistor-end1) % Verbinung zu Wiederständen 
                node[circle, fill, inner sep=1.5pt]{}; % Punkt am Ende der Verbindung 
            \draw (3,-1) node[nmos, rotate=270] (nmosR2) {}
                node at (nmosR2.center) [anchor=center, yshift=-10pt] {$N_4$}
                (nmosR2.source) -- (resistor-end2)
                node[circle, fill, inner sep=1.5pt]{};

            % NMOS Drain Verbindung
            \draw (nmosR1.drain) -- ++ (0.25,0) % horizontal nach rechts
                -- ++ (0,-3) coordinate (comp) % vertikale Linie 
                node[circle, fill, inner sep=1.5pt]{};

            \draw (nmosR1.gate) -- ++(0,0.25) -- ++(-0.75,0) 
                node[above, midway]{$V_{res}$};
            \draw (nmosR2.gate) -- ++(0,0.25) -- ++(-0.75,0)
                node[above, midway]{$V_{res0}$};

            % Dreieck am Ende von Vcomp
            \draw (6.75,-0.51) node[op amp] (triangle){}
                (triangle.+) -- (nmosR2.drain);
            \draw (triangle) node[above, yshift=27pt] {Komparator};  

            \draw (triangle.-) -- ++(-0.5,0) node[above]{$V_{ref}$};

            % Verbindung Ground
            \draw (nmosL.source) |- (resistor-end3)
                (nmosL2.source) |- (resistor-end3) coordinate[pos=0.5](gndM)
                (gndM) -- ++(0,-0.25) node[sground]{}
                    node[right, xshift=5pt, yshift=-5pt]{Gnd}
                (gndM) node[circle, fill, inner sep=1.5pt]{}; 

            % Verbindung VDD
            \draw (resistor-begin) |- (resistor-begin1)
                (nmosL1.drain) |- (resistor-begin1) coordinate[pos=0.5](vddM)
                (vddM) node[circle, fill, inner sep=1.5pt]{}
                (vddM) -- ++(0,0.5) coordinate(vddA)
                (vddA) -- ++(-0.5,0)
                (vddA) -- ++(0.5,0)
                (vddA) node[above] {$V_{DD}$};

        \end{circuitikz}
    \end{center}
    \caption[Generierung des Referenzstrom durch Transistor]{\\Quelle: Generierung des Referenzstrom durch Transistor (Eigene Darstellung)}
    \label{fig:Vref&Vres}
\end{figure}
%\vspace{1.5\baselineskip}
Bei dem \gls{nmos}-Transistor $N_1$ handelt es sich um einen Native Transistor, da es sonst zu einer Verzögerung in der Generierung der Referenzspannung kommt und der Schaltvorgang des Komparator nicht richtig ausgeführt werden kann. Der Transistor $N_2$ fungiert als \gls{nmos}-Diode. Die so entstehende Referenzspannung $V_{ref}$ hängt von der Wahl des Widerstands $R_0$ sowie von den transistorabhängigen Parametern wie der Threshold-Spannung ($V_{th}$) und dem Stromverstärkungsfaktor ab. Die Schaltung stellt sicher, dass $V_{ref}$ über einen weiten Bereich stabil bleibt und unabhängig von Versorgungsspannungsschwankungen oder externen Störungen eine konstante Referenz bildet \cite{CircuitDesigLayoutAndSimulation}. Durch diese Entwurfsanpassung konnte die Schwankung der Referenzspannung um ca. $^1/_3$ auf $\approx$ 99 mV wie in \autoref{fig:Vref_Variation} zu sehen ist, verringert werden. Dies führt dazu, dass die Schaltung auch über die Corner-Simualation den Spezifikationen entspricht. \autoref{tab:corner_summary1.5} zeigt die wichtigsten Analysewerte der Simulation für die \gls{por}-Schaltung des digitalen Teils: 
\begin{table}[H]
    \centering
    \small
    \begin{tabular}{c|l|c|c|c}
        {\rotatebox{90}{\parbox{1.8cm}{\textbf{Spannung}}}} &
        \textbf{\makecell{Werte}} &
        \textbf{\makecell{Spezifikation}} &
        \textbf{\makecell{Min. Wert}} &
        \textbf{\makecell{Max. Wert}} \\
        \midrule

        \multirow{5}{*}{\rotatebox{90}{5 V}} &
        Verzögerung steigende Flanke & 20 us $\leq$ x $\leq$ 155 us & 20.78 us & 153.9 us \\
        & Verzögerung fallende Flanke & $\leq$ 1 us & 48.99 ns & 206.3 ns \\
        & Schwellspannung steigend & 1.1 V $\leq$ x $\leq$ 1.45 V & 1.133 V & 1.443 V \\
        & Schwellspannung fallend & 1.05 V $\leq$ x $\leq$ 1.4 V & 1.066 V & 1.359 V \\
        & Hysterese & 50 mV $\leq$ x $\leq$ 90 mV & 66.97 mV & 84.79 mV \\
        \midrule
    \end{tabular}
    \caption{Ergebnisse Corner-Simulation PoR-Schaltung für den digitalen Teil}
    \label{tab:corner_summary1.5}
\end{table}





































Die Analysewerte der Corner-Simulation der \gls{por}-Schaltung für den analogen Schaltungsteil werden in \autoref{tab:corner_summary} dargestellt. Dabei überschreitet der Maximalwert für den fallenden Schwellenwert mit 15 mV und befindet sich damit außerhalb der Spezifikation. Da es sich bei der Corner-Simulation um ein Worst-Case-Szenario handelt, ist bei dieser geringen Abweichung keine weitere Anpassung notwendig.
\begin{table}[H]
    \centering
    \small
    \begin{tabular}{c|l|c|c|c}
        {\rotatebox{90}{\parbox{1.8cm}{\textbf{Spannung}}}} &
        \textbf{\makecell{Werte}} &
        \textbf{\makecell{Spezifikation}} &
        \textbf{\makecell{Min. Wert}} &
        \textbf{\makecell{Max. Wert}} \\
        \midrule

        \multirow{5}{*}{\rotatebox{90}{5 V}} &
        Verzögerung steigende Flanke & 20 us $\leq$ x $\leq$ 130 us & 24.37 us & 126.5 us \\
        & Verzögerung fallende Flanke & $\leq$ 1 us & 34.53 ns & 138 ns \\
        & Schwellspannung steigend & 2.2 V $\leq$ x $\leq$ 2.6 V & 2.206 V & 2.577 V \\
        & Schwellspannung fallend & 2.1 V $\leq$ x $\leq$ 2.5 V & 2.153 V & \cellcolor{lightyellow}2.515 V \\
        & Hysterese & 50 mV $\leq$ x $\leq$ 120 mV & 52.96 mV & 63.08 mV \\
        \midrule
    \end{tabular}
    \caption{Ergebnisse Corner-Simulation PoR-Schaltung für den analogen Teil}
    \label{tab:corner_summary}
\end{table}

\section{Monte-Carlo-Simulation}
Nachdem die \gls{por}-Schaltung die Spezifikationen in der Corner Simulation erfüllt, wurde abschließend eine \gls{mc}-Simulation durchgeführt. Diese Simulation testet die Zuverlässigkeit der Schaltung unter realen Fertigungsbedingungen. Während die Corner-Simulation systematische Schwankungen (Worst-Case-Szenarien) untersucht, ermöglicht die \gls{mc}-Simulation eine statistische Analyse zufälliger Variationen, die während der Herstellung auftreten können. In Cadence können \gls{mc}-Simulationen in drei verschiedene Kategorien unterteilt werden, welche in diesem Fall mit jeweils 200 Iterationen ausgeführt wurden. Das bedeutet, für jeden Durchlauf wurden den Bauteilen zufällige Parameter (in Form einer Gaußverteilung), die durch die Fertigung auftreten können, zugeteilt und berechnet. 

Bei der Process-Simulation innerhalb der \gls{mc}-Simulation werden alle Transistoren und Bauteile mit zufälligen Abweichungen versehen, welche globale Fertigungsschwankungen widerspiegeln. Diese Schwankungen wirken sich auf alle Bauteile einer Wafer-Charge\footnote{\url{https://heinen-elektronik.de/glossar/wafer/} --\,zuletzt am 27.~Februar~2025 erfolgreich aufgerufen.} gleichermaßen aus. 

Die Mismatch-Simulation untersucht lokale Unterschiede zwischen identischen Bauelementen. Dazu zählen beispielsweise kleine Abweichungen zwischen den Transistoren, die einen Differenzverstärker bilden. Das Ergebnis ist eine Verteilung der Offset-Spannung oder Stromabweichung aufgrund von Fertigungsunterschieden zwischen benachbarten Bauelementen.

Die All-Simulation ist eine Kombination aus lokalen Mismatch-Variationen und globalen Prozessabweichungen, um ein realitätsnahes Verhalten der Schaltung zu simulieren. Sie stellt sicher, dass weder großflächige noch feinere Variationen in der Fertigung sich außerhalb der spezifizierten Grenzen befinden. 

\subsection{Einfluss von Flankenzeiten auf die Funktionalität}
Die drei Varianten der \gls{mc}-Simulation werden unter Variation der Anstiegs- und Abfallzeiten (Flanken) der Versorgungsspannung ausgeführt. Die Anstiegszeit beschreibt die Zeit, welche die Versorgungsspannung benötigt, um von 0~V auf 1.5~V (für den digital Teil) oder 5~V~(für den analogen Teil) anzusteigen. Bei der Abfallzeit handelt es sich dementsprechend um die Dauer des Abfallen der Versorgungsspannung beim Ausschalten des Systems. Diese Parameter sind besonders für \gls{por}-Schaltungen von großer Bedeutung, da sie direkt beeinflussen, wann die interne Referenzspannung aufgebaut wird und wie schnell der Komparator in den aktiven Zustand übergeht. 

Bei kurzen Anstiegszeiten (z.B. 10 ms) steigt die Versorgungsspannung sehr schnell an, wodurch sich die Schaltung innerhalb kürzester Zeit initialisieren muss. Dies stellt erhöhte Anforderungen an die Generierung der Referenzspannung und an die Stabilität des Komparators bei schnell wechselnden Flanken. Besonders kritisch ist dabei die zuverlässige Erkennung des Spannungsniveaus sowie die saubere Generierung des Reset-Signals, bevor nachfolgende Schaltungsteile aktiviert werden. Längere Anstiegszeiten (z.B. 100 ms) führen hingegen zu einem langsameren Spannungsanstieg. Dabei hat die Schaltung mehr Zeit, die interne Referenzspannung aufzubauen und die Eingangsbedingungen des Komparators zu stabilisieren. Allerdings können bei sehr langsamen Anstiegszeiten unerwünschte Zwischenzustände auftreten, bei denen die Spannung bereits hoch genug ist, um Teilschaltungen zu aktivieren, während die Reset-Bedingungen noch nicht sicher erfüllt sind. Dies kann insbesondere in Kombination mit Prozess- und Mismatch-Schwankungen zu einer erhöhten Fehleranfälligkeit führen.

Ähnliche Effekte treten bei der Variation der fallenden Flanke auf. Ein schnelles Abschalten~(10 ms) entlädt die Versorgung sehr abrupt, was dazu führen kann, dass das Reset-Signal nicht rechtzeitig ausgelöst wird. Bei einem langsamen Spannungsabfall (100 ms) können unerwünschten Restströmen oder Fehlinterpretationen von Schaltzuständen auftreten. Damit sich keine falschen Zustände im System einstellen, muss das Reset-Signal sowohl für langsame als auch schnell fallende Flanken zuverlässig ausgelöst werden. 

\subsection{Analyse der Abweichungen in der Monte-Carlo-Simulation}
Die \gls{mc}-Process-Simulation zeigte, dass die Schwellwerte der steigenden und fallenden Flanken fast 200 mV außerhalb der spezifizierten Grenzen liegen wie in \autoref{tab:MCProcess100} zu sehen ist. Dieses Verhalten ist insofern ungewöhnlich, da die Corner-Simulationen, die die Worst-Case-Bedingungen abdecken sollen. Die in der Spezifikation angegeben Werte werden an den geprüften Temperaturpunkten (-40 $^{\circ}$C, 27 $^{\circ}$C und 125 $^{\circ}$C) sicher einhalten. Dass die \gls{mc}-Process-Simulation bei nominaler Temperatur von 27 $^{\circ}$C eine so deutliche Abweichung aufweist, deutet auf eine Unstimmigkeit im Simulations- oder Modellierungsprozess hin. Aus diesem Grund muss die Dokumentation der Cadence-Simulation noch einmal genauer überprüft, um die Ursache dieser Abweichung herauszufinden. Bisherige Tests deuten darauf hin, dass der Fehler durch den nativ Transistor verursacht wird, auch wenn die Wert der \gls{por}-Schaltung für den digitalen Teil den Spezifikationen entsprechen. 
\begin{table}[H]
    \centering
    \small
    \begin{tabular}{c|l|c|c|c}
        {\rotatebox{90}{\parbox{1.8cm}{\textbf{Spannung}}}} &
        \textbf{\makecell{Werte}} &
        \textbf{\makecell{Spezifikation}} &
        \textbf{\makecell{Min. Wert}} &
        \textbf{\makecell{Max. Wert}} \\
        \midrule

        \multirow{5}{*}{\rotatebox{90}{5 V}} &
        Verzögerung steigende Flanke & 20 us $\leq$ x $\leq$ 130 us & 46.25 us & 89.11 us \\
        & Verzögerung fallende Flanke & x $\leq$ 1 us & 57.65 ns & 79.79 ns \\
        & Schwellspannung steigend & 2.2 V $\leq$ x $\leq$ 2.6 V & \cellcolor{lightred}2.088 V & \cellcolor{lightred}2.722 V \\
        & Schwellspannung fallend & 2.1 V $\leq$ x $\leq$ 2.5 V & \cellcolor{lightred}2.038 V & \cellcolor{lightred}2.659 V \\
        & Hysterese & 50 mV $\leq$ x $\leq$ 120 mV & \cellcolor{lightyellow}49.92 mV & 62.88 mV \\
        \midrule

        \multirow{5}{*}{\rotatebox{90}{1.5 V}} &
        Verzögerung steigende Flanke & 20 us $\leq$ x $\leq$ 230 us & 42.36 us & 77.56 us \\
        & Verzögerung fallende Flanke & x $\leq$ 1 us & 80.05 ns & 118.6 ns \\
        & Schwellspannung steigend & 1.1 V $\leq$ x $\leq$ 1.45 V & 1.147 V & 1.445 V \\
        & Schwellspannung fallend & 1.05 V $\leq$ x $\leq$ 1.4 V & 1.079 V & 1.361 V \\
        & Hysterese & 50 mV $\leq$ x $\leq$ 120 mV & 68.14 mV & 84.18 mV \\
        \midrule
        
    \end{tabular}
    \caption{MC-Process-Simulation mit Anstiegs-und Abfallzeit von 100~ms}
    \label{tab:MCProcess100}
\end{table}

Bei den in \autoref{tab:MCMismatch100} zu sehenden Ergebnissen der \gls{mc}-Mismatch-Simulation liegen die Schwellwerte der Flanken und die anderen entschiedenen Messwerte innerhalb der Spezifikation. Auch die Schwankungsbreite der Ergebnisse ist bei der Mismatch-Simulation geringer als bei der Corner-Simulation, was grundsätzlich dem erwarteten Verhalten entspricht. Diese Beobachtung gilt sowohl für die \gls{por}-Schaltung im analogen (5~V) als auch im digitalen (1,5~V) Schaltungsteil.
\begin{table}[H]
    \centering
    \small
    \begin{tabular}{c|l|c|c|c}
        {\rotatebox{90}{\parbox{1.8cm}{\textbf{Spannung}}}} &
        \textbf{\makecell{Werte}} &
        \textbf{\makecell{Spezifikation}} &
        \textbf{\makecell{Min. Wert}} &
        \textbf{\makecell{Max. Wert}} \\
        \midrule
        
        \multirow{5}{*}{\rotatebox{90}{5 V}} &
        Verzögerung steigende Flanke & 20 us $\leq$ x $\leq$ 130 us & 40.02 us & 108.9 us \\
        & Verzögerung fallende Flanke & x $\leq$ 1 us & 63.88 ns & 72.03 ns \\
        & Schwellspannung steigend & 2.2 V $\leq$ x $\leq$ 2.6 V & 2.35 V & 2.457 V \\
        & Schwellspannung fallend & 2.1 V $\leq$ x $\leq$ 2.5 V & 2.293 V & 2.4 V \\
        & Hysterese & 50 mV $\leq$ x $\leq$ 120 mV & 45.44 mV & 57.49 mV \\
        \midrule

        \multirow{5}{*}{\rotatebox{90}{1.5 V}} &
        Verzögerung steigende Flanke & 20 us $\leq$ x $\leq$ 230 us & 33.84 us & 77.86 us \\
        & Verzögerung fallende Flanke & x $\leq$ 1 us & 84.6 ns & 108 ns \\
        & Schwellspannung steigend & 1.1 V $\leq$ x $\leq$ 1.45 V & 1.219 V & 1.343 V \\
        & Schwellspannung fallend & 1.05 V $\leq$ x $\leq$ 1.4 V & 1.146 V & 1.265 V \\
        & Hysterese & 50 mV $\leq$ x $\leq$ 120 mV & 72.77 mV & 78.94 mV \\
        \midrule
        
    \end{tabular}
    \caption{MC-Mismatch-Simulation mit Anstiegs-und Abfallzeit von 100~ms}
    \label{tab:MCMismatch100}
\end{table}

Da die \gls{mc}-All-Simulation eine Kombination aus Process- und Mismatch-Variatio\-nen ist, treten die beschriebenen Abweichungen auch bei dieser Simulation auf (siehe \autoref{tab:MCAll100}). Die Schwellwerte der Flanken liegen dabei ebenfalls außerhalb der spezifizierten Grenzen. Bei dieser Simulation sind auch die Wert der \gls{por}-Schaltung bei einer Betriebsspannung von 1.5~V außerhalb der Werte der Corner-Simulation. 
\begin{table}[H]
    \centering
    \small
    \begin{tabular}{c|l|c|c|c}
        {\rotatebox{90}{\parbox{1.8cm}{\textbf{Spannung}}}} &
        \textbf{\makecell{Werte}} &
        \textbf{\makecell{Spezifikation}} &
        \textbf{\makecell{Min. Wert}} &
        \textbf{\makecell{Max. Wert}} \\
        \midrule

        \multirow{5}{*}{\rotatebox{90}{5 V}} &
        Verzögerung steigende Flanke & 20 us $\leq$ x $\leq$ 130 us & 38.94 us & 115.5 us \\
        & Verzögerung fallende Flanke & x $\leq$ 1 us & 57.15 ns & 81.63 ns \\
        & Schwellspannung steigend & 2.2 V $\leq$ x $\leq$ 2.6 V & \cellcolor{lightred}2.058 V & \cellcolor{lightred}2.75 V \\
        & Schwellspannung fallend & 2.1 V $\leq$ x $\leq$ 2.5 V & \cellcolor{lightred}2.008 V & \cellcolor{lightred}2.686 V \\
        & Hysterese & 50 mV $\leq$ x $\leq$ 120 mV & \cellcolor{lightyellow}49.47 mV & 64.14 mV \\
        \midrule

        \multirow{5}{*}{\rotatebox{90}{1.5 V}} &
        Verzögerung steigende Flanke & 20 us $\leq$ x $\leq$ 230 us & 36.93 us & 88.91 us \\
        & Verzögerung fallende Flanke & x $\leq$ 1 us & 77.38 ns & 129.2 ns \\
        & Schwellspannung steigend & 1.1 V $\leq$ x $\leq$ 1.45 V & 1.123 V & \cellcolor{lightyellow}1.452 V \\
        & Schwellspannung fallend & 1.05 V $\leq$ x $\leq$ 1.4 V & 1.055 V & 1.368 V \\
        & Hysterese & 50 mV $\leq$ x $\leq$ 120 mV & 67.06 mV & 83.75 mV \\
        \midrule
        
    \end{tabular}
    \caption{MC-All-Simulation mit Anstiegs-und Abfallzeit von 100~ms}
    \label{tab:MCAll100}
\end{table}

Weitere Analysen konzentrieren sich daher speziell auf die Modellparameter des nativen Transistors sowie die genaue Prozessbeschreibung innerhalb der Cadence-Umgebung, um die Abweichung einzugrenzen und gegebenenfalls Korrekturmaßnahmen einzuleiten oder den Entwurf noch mal anzupassen. 

Die Ergebnisse der \gls{mc}-Simulationen bei Anstiegs-und Abfallzeiten von 10~ms~(siehe~\ref{tab:MCProcess10}, \ref{tab:MCMismatch10}, \ref{tab:MCAll10}) liegen unterscheiden sich nur gering von den Ergebnissen der Simulationen bei 100~ms. In beiden Fällen bewegen sich die Werte, abgesehen von den bekannten Abweichungen durch den Nativ Transistor, innerhalb der spezifizierten Grenzen.  

Um sicherstellen zu können, dass das Reset-Signal auch bei schnellem Flankenwechsel zuverlässig funktioniert, wurde noch eine \gls{mc}-All-Simulation mit einer Anstiegs-und Abstiegszeit von nur 1~ms durchgeführt. Wie zuvor liegen auch bei dieser Simulation die Werte der Schwellspannung außerhalb der Spezifikation, dennoch bleibt die Grundlegende Funktion der Schaltung erhalten. Die Verzögerung und Freigabe des Reset-Signals findet innerhalb der spezifizierten Grenzen statt (siehe \ref{tab:MCAll1}), wodurch eine zuverlässige Funktion der \gls{por}-Schaltungen bestätigt werden konnte. 